% ============================================================
% DRAFT NOTE (intervention)
% Contains: The parity ladder that validates the causal interpretation and
% documents when live forks, re-prefill, or historical labels are admissible.
% Written now: The full validation requirements and their rationale.
% Pending: Per-compressor and per-ratio parity outcomes.
% Sources: docs/methodology.md; docs/_why/6_intervention_semantics.md;
% docs/goal.md
% ============================================================
\section{Validated Intervention Protocol}
\label{sec:intervention}

The paired estimand is only meaningful when the two branches differ solely in
the intended compression operation. We therefore treat intervention validation
as a first-class contribution rather than an implementation detail. Each
compressor is classified as cache-native, auxiliary-state, or
re-prefill-defined before its labels are interpreted. Live-state forks are
preferred for the first two classes; matched re-prefill is the faithful target
for the third.

The validation ladder has three checks. First, fork-clone parity and isolation
continue two uncompressed state forks and require identical token IDs, while
also verifying that writable cache storage, execution order, and hooks cannot
cross-contaminate the branches. Second, sham re-prefill parity compares a live
full-cache continuation with a no-press re-prefill from the same prefix.
Third, where a live operator exists, intervention parity compares the
compressed live-state fork with pressed re-prefill from that prefix. Every
check records exact continuation match rate, first divergence position, and
final quality delta.

Passing sham parity permits a stored continuous reference to stand in for a
no-press re-prefill. Passing intervention parity permits historical pressed
re-prefill hybrids to be described as equivalent to live-cache activation. If
either condition fails, labels are generated under the corresponding matched
target instead of relying on a shortcut.

\begin{table}[t]
\centering
\caption{Parity outcomes, reported for every compressor and removal ratio.}
\label{tab:parity}
\begin{tabular}{lllll}
\toprule
Compressor & Ratio & Fork and isolation & Sham re-prefill & Intervention \\
\midrule
\todo{compressor} & \todo{ratio} & \todo{outcome} & \todo{outcome} & \todo{outcome} \\
\todo{compressor} & \todo{ratio} & \todo{outcome} & \todo{outcome} & \todo{outcome} \\
\bottomrule
\end{tabular}
\end{table}

Figure~\todo{F2 reference} will depict the ladder and the branch pairs it
validates.
